\section{Geometry and Space}
\label{sec:geometry}

There is no container. Space is a derived description of the note-structure. Distance is how far apart two vibes are in the web of notes, and the large-scale shape of that web is the geometry of the world.

\subsection{Space is note-distance}

Two vibes joined by strong notes are close. Weakly joined, they are far. The fundamental metric is the shortest path through the note-structure, weighted so that strong notes are short. This is purely combinatorial, hop counts and integer strengths rather than real coordinates. Closeness is literally strong notes. Space is made of the relations, not a stage they sit on, and smooth geometry is the large-scale limit only. By the principle of no hidden state (Section~\ref{sec:constraints}), a vibe holds no coordinates. Its position is its place in the note-structure, computable from local notes by a finite procedure, not a label it carries.

\subsection{What geometry, and why}

If the mesh is a regular tiling $\{p, q\}$, meaning $p$-sided cells with $q$ meeting at each corner, one integer inequality fixes its geometry,
\[
(p-2)(q-2) < 4 \ \text{spherical}, \quad (p-2)(q-2) = 4 \ \text{flat}, \quad (p-2)(q-2) > 4 \ \text{hyperbolic}.
\]
So the question of what geometry the world has reduces to two integers. The model leans hyperbolic, where $(p-2)(q-2) > 4$, because hyperbolic space gives what the framework wants and the flat plane cannot. A vibe can have a small set of immediate notes yet be only a few notes away from vast numbers of other vibes, since the number of vibes within $n$ notes grows exponentially. This matches the eternal-expansion commitment of Section~\ref{sec:emergence}. The substrate has unbounded room a few steps out.

\begin{proposition}[The exponential tree]
Take a regular tree in which every node has $q$ branches, a Bethe lattice. Each node touches only $q$ neighbors, yet the number of nodes within $n$ steps grows like $q^n$, exponentially. This is hyperbolic geometry in its barest form, small local degree with vast reach a few steps out. Note-distance on such a tree is just the number of branches between two nodes, distance made entirely of connections.
\end{proposition}

A hyperbolic mesh is a tree wearing a geometric costume, so the nesting hierarchy of selves (Section~\ref{sec:mind}) and the spatial layout are the same object viewed two ways.

\subsection{Curvature, gravity, and dimension}

Curvature is computed from the note-combinatorics. Positive curvature marks a tightly bound, mesh-like region, a coherent self. Negative curvature marks a sparsely bridged, tree-like region. Gravity is the tendency of vibes to densify their notes and reduce isolation, a flow toward more positive curvature where coherent structures form. Mass concentrates curvature, the discrete echo of mass curving space. A striking discrete bonus is that regular hyperbolic honeycombs exist only up to four dimensions, with nothing in five. If the world wants exponential reach together with regularity, the geometry caps at four dimensions, a candidate, cost-free reason a world of three space dimensions and one of time is special.

\subsection{Spacetime}

Space alone is the note-metric. Spacetime is that metric together with the causal order of Section~\ref{sec:time}. The partial order of which update could influence which supplies the light-cone structure, and the note-distance supplies the spatial part. Together they are a causal set, the discrete-spacetime object. The one-note-per-beat bound is the discrete speed of light, and differential beat rates give the time-dilation analogue.
